% !PRINTED 205
\weberpage{205}
%<<CONT
\gesperrt{transform the group so that it contains the two
substitutions}
\[
\Theta = \varepsilon x, \quad \psi = \frac{c}{x}
\]
\gesperrt{where $\varepsilon$ denotes a primitive $\nu$th root of unity,
and $c$ can be an arbitrarily prescribed
value, e.g. also $1$.}

The hypothesis of this theorem is, in the cases enumerated
in \S.\,52, always satisfied for one of the various values $\nu$,
except in the case of the cyclic group and of the
dihedral group with $m = 2$.

Finally we prove also the following theorem:

2.~\gesperrt{If $a, a'$ are the poles of a substitution $\Theta$ of $G$,
then the poles conjugate to $a$ and $a'$,}
\[
b = \chi(a), \quad b' = \chi(a'),
\]
\gesperrt{where $\chi$ is an arbitrary substitution of $G$, are the
two poles of one and the same substitution,
namely of the substitution $\chi\Theta\chi^{-1}$.}

The correctness follows immediately from inspection
of the equations:
\[
\begin{aligned}
a &= \Theta(a), & \chi\Theta\chi^{-1}\chi(a) &= \chi\Theta(a) = \chi(a) \\
a' &= \Theta(a'), & \chi\Theta\chi^{-1}\chi(a') &= \chi\Theta(a') = \chi(a').
\end{aligned}
\]

\begin{center}
\S.\,54.

\gesperrt{The Fundamental Forms.}
\end{center}

In order to arrive at the definitive enumeration of all finite groups
$G$, it is necessary to enter more closely into the invariants of the
corresponding binary substitution groups.

We must therefore, besides the linear fractional substitutions


\weq{1}{y = \frac{\alpha x + \beta}{\gamma x + \delta},}

which we have denoted by $A'''$ in \S.\,49, also the binary
substitutions

\weq{2}{(y_1, y_2) = \begin{pmatrix}\alpha, & \beta \\ \gamma, & \delta\end{pmatrix}(x_1, x_2)}
%CONT>>
